\chapter{Constructing the rotational bowl soliton}\label{sec:examples-bowl}
The radial equation in Proposition \ref{prop:radial-translator-tip} admits a complete solution, not just a formal series at the tip. We give the construction and its leading asymptotics in separate steps. Compare \cite{CSS07}, Section 2, equations (2.1)--(2.2) and Lemma 2.2; citations use the author-hosted September 2005 manuscript. Our local construction at the singular endpoint is supplied explicitly. Throughout this section $n\geq2$ and the translation speed is one upward.

\begin{lemma}[Resolving the radial equation at zero]\label{lem:bowl-smooth-tip-existence}
There is $R>0$ and a unique smooth radial solution $u(x)=f(|x|)$ on $B_R\subset\mathbb R^n$ of the unit-speed translator equation with $f(0)=0$. Its derivative has the form
\[
 p(r)=f'(r)=rq(r^2),\qquad q(0)=\frac1n,
\]
where $q$ is real analytic near zero. In particular the apparent $1/r$ coefficient in the radial equation creates no singularity in the graph at its tip.
\end{lemma}
\begin{proof}
Put $m=n-1$ and $P(q)=q^2(1-mq)$. Multiplying $p'+mp/r=1+p^2-mp^3/r$ by $r^{n-1}$ and integrating from zero gives, after setting $p=rq(r^2)$,
\[
 q(z)=\frac1n+z\int_0^1t^{n+1}P(q(zt^2))\,dt.
\]
To construct a smooth solution without assuming it exists, use the Banach space of functions continuous on a small closed complex disk $|z|\leq\delta$ and holomorphic in its interior, with the supremum norm. On the closed ball $\|q-1/n\|_\infty\leq1$, the polynomial $P$ is bounded by some $M_n$ and is Lipschitz with some $L_n$. The right side of (*) maps this ball into itself when $\delta M_n/(n+2)\leq1$ and is a contraction when $\delta L_n/(n+2)<1$. Its unique fixed point is holomorphic and real on the real axis, since iteration starting from $1/n$ preserves complex conjugation. Define
\[
 f(r)=\frac12\int_0^{r^2}q(s)\,ds.
\]
This is analytic as a function of $|x|^2$ and satisfies the ODE away from zero by differentiation of the integral identity. At zero its Hessian is $I/n$, so the graph equation holds there as well. Conversely any smooth radial solution has $p/r\to1/n$ by Proposition \ref{prop:radial-translator-tip}, and it obeys the same integral equation on the real interval. The same contraction proof in the space of continuous real functions gives local uniqueness without assuming analyticity of the competitor. The only external analytic tools in this local construction are the contraction principle and uniform-limit closure of holomorphic functions.
\end{proof}

\begin{lemma}[Linear barriers prevent finite-radius blow-up]\label{lem:bowl-slope-barriers}
The tip solution extends to all $r\geq0$ and satisfies
\[
 \frac rn<p(r)<\frac r{n-1},\qquad p'(r)>0\quad(r>0).
\]
Consequently its height obeys
\[
 \frac{r^2}{2n}<f(r)<\frac{r^2}{2(n-1)}\quad(r>0).
\]
\end{lemma}
\begin{proof}
The tip expansion in Proposition \ref{prop:radial-translator-tip} gives $p=r/n+r^3/(n^3(n+2))+o(r^3)$, so the strict barriers hold initially. At a contact with $p=r/n$, the ODE gives $p'-1/n=r^2/n^3>0$, preventing a first downward crossing. At a contact with $p=r/(n-1)$, it gives $p'-1/(n-1)=-1/(n-1)<0$, preventing a first upward crossing. Thus the two barriers persist for the whole interval of existence. Between them the factor $1-(n-1)p/r$ is positive, so $p'>0$. If a finite positive endpoint of existence existed, the barriers would bound $p$ there, and the smooth ODE away from $r=0$ would continue the solution across that endpoint. This contradiction proves global extension. Integrating the strict bounds from zero gives the height bounds. Compare the continuation and barrier argument in \cite{CSS07}, proof of Lemma 2.1.
\end{proof}

\begin{theorem}[The complete rotational bowl]\label{thm:rotational-bowl-construction}
For every $n\geq2$, the graph of the function $f(|x|)$ constructed above is a smooth, properly embedded, complete, strictly convex hypersurface in $\mathbb R^{n+1}$ translating with velocity $e_{n+1}$. It is unique among smooth rotational entire translating graphs with this speed, rotation axis, and tip height zero.
\end{theorem}
\begin{proof}
Lemma \ref{lem:bowl-smooth-tip-existence} gives smoothness at the origin, and Lemma \ref{lem:bowl-slope-barriers} gives global continuation. The Hessian eigenvalues for $r>0$ are $p'>0$ and $p/r>0$; at zero they are all $1/n$. Thus the graph is strictly convex. The ODE and Proposition \ref{prop:translator-graph-example} give its translating equation. An entire smooth graph is embedded and proper: the base projection bounds the preimage of an ambient compact set, which is also closed. Its induced metric dominates the Euclidean base metric, proving completeness by the Cauchy-sequence argument used for the catenoid. Local uniqueness at the tip and ordinary ODE uniqueness on $r>0$ give uniqueness in the stated rotational class. This is the rotational existence result in \cite{CSS07}, Lemma 2.2; no classification of nonrotational translators is claimed.
\end{proof}

\begin{proposition}[The quadratic growth coefficient of the bowl]\label{prop:bowl-leading-growth}
For the normalized bowl,
\[
 \lim_{r\to\infty}\frac{p(r)}r=\frac1{n-1},\qquad
 \lim_{r\to\infty}\frac{f(r)}{r^2}=\frac1{2(n-1)}.
\]
\end{proposition}
\begin{proof}
Fix $b$ with $1/n<b<1/(n-1)$. If $p(r)<br$ for every sufficiently large $r$, the ODE and $p\geq r/n$ imply $p'\geq(1-(n-1)b)r^2/n^2$. Integration would force cubic growth, contradicting the linear upper bound in Lemma \ref{lem:bowl-slope-barriers}. Therefore $p$ reaches $br$ at some arbitrarily large radius. At such a contact the derivative difference is
\[
 p'-b=(1+b^2r^2)(1-(n-1)b)-b>0
\]
for large $r$, so it cannot subsequently cross downward. Thus $\liminf p/r\geq b$. Let $b\uparrow1/(n-1)$ and use the upper barrier to get the first limit. Integrating $(1/(n-1)-\varepsilon)r\leq p(r)\leq r/(n-1)$ outside a fixed radius, then dividing by $r^2$, proves the second. This is the first asymptotic barrier step in \cite{CSS07}, Lemma 2.1.
\end{proof}

\begin{lemma}[The first slope correction]\label{lem:bowl-slope-correction}
Write $m=n-1$ and $w(r)=r/m-p(r)>0$. Then
\[
 \lim_{r\to\infty}rw(r)=1.
\]
\end{lemma}
\begin{proof}
The equation for the defect is
\[
 w'=\frac1m-a(r)w,\qquad a(r)=\frac m r(1+p(r)^2),
 \qquad \frac{a(r)}r\longrightarrow\frac1m.
\]
Fix $0<\varepsilon<1$ and compare $w$ with $b_\pm=(1\pm\varepsilon)/r$. At $w=b_+$ the right side tends to $-\varepsilon/m$, while $b_+'\to0$; for large radii $b_+$ is a strict supersolution. If $w>b_+$ persisted forever, then $w'\leq-\varepsilon/(2m)$ for all sufficiently large radii, contradicting $w>0$. Thus $w$ reaches this upper barrier and cannot cross it upward afterward. Similarly $b_-$ is a strict subsolution for large radii. If $w<b_-$ persisted forever, then $w'\geq\varepsilon/(2m)$, contradicting $w<b_-\to0$. Once it crosses upward it cannot cross back. Therefore $1-\varepsilon\leq rw\leq1+\varepsilon$ eventually. Let $\varepsilon\downarrow0$.
\end{proof}

\begin{proposition}[A logarithmic correction to the bowl height]\label{prop:bowl-logarithmic-asymptotic}
There is a constant $C_n$ determined by the normalization $f(0)=0$ such that
\[
 p(r)=\frac r{n-1}-\frac1r+O(r^{-3}),\qquad
 f(r)=\frac{r^2}{2(n-1)}-\log r+C_n+O(r^{-2}).
\]
The constant cannot be discarded unless one changes the vertical translation of the bowl.
\end{proposition}
\begin{proof}
With $m=n-1$, the exact equation for $w$ from Lemma \ref{lem:bowl-slope-correction} is $w'=F(r,w)$, where
\[
 F(r,w)=\frac1m-\left(\frac r m+\frac m r\right)w+2w^2-\frac m r w^3.
\]
For $b_K(r)=r^{-1}+Kr^{-3}$, expansion gives
\[
 F(r,b_K)-b_K'
 =\frac{3-m-K/m}{r^2}+O_K(r^{-4}).
\]
Choose $K_+>m(3-m)$ and $K_-<m(3-m)$. Then $b_{K_+}$ is a supersolution and $b_{K_-}$ a subsolution for $r\geq R$ large. By Lemma \ref{lem:bowl-slope-correction}, both $w$ and these barriers lie between $0$ and $2/r$ there. Throughout this interval
\[
 \partial_wF=-\frac r m-\frac m r+4w-\frac{3m}r w^2
 \leq-\frac r{2m}
\]
after increasing $R$. On any interval where $w-b_{K_+}>0$, the mean-value theorem gives $(w-b_{K_+})'\leq-r(w-b_{K_+})/(2m)$. At a zero this difference cannot cross from negative to positive. Hence its positive part is at most $A\exp(-(r^2-R^2)/(4m))$. Applying the same argument to $(b_{K_-}-w)_+$ gives the lower estimate. The exponentially decaying errors can be absorbed in $O(r^{-3})$, so $w=1/r+O(r^{-3})$. This proves the slope expansion. The derivative of $f-r^2/(2m)+\log r$ is consequently $O(r^{-3})$, which is integrable at infinity; its tail integral is $O(r^{-2})$. This proves the height expansion, including its additive constant. Compare \cite{CSS07}, Lemmas 2.1--2.2; the argument here uses scalar barriers rather than the longer formal series.
\end{proof}

\begin{corollary}[Translation speed and dilations of the bowl]\label{cor:bowl-speed-scaling}
For any $c>0$, define $f_c(r)=c^{-1}f(cr)$. Its graph translates with velocity $ce_{n+1}$ and has tip Hessian $(c/n)I$. Its large-radius expansion is
\[
 f_c(r)=\frac{cr^2}{2(n-1)}-\frac1c\log r
 +\frac{C_n-\log c}{c}+O_c(r^{-2}).
\]
\end{corollary}
\begin{proof}
We have $f_c'=p(cr)$ and $f_c''=c p'(cr)$. Substitution in the radial graph operator multiplies its unit-speed value by $c$. Smoothness, completeness and properness follow by dilation. At zero the second derivative is $c/n$. Replacing $r$ by $cr$ in Proposition \ref{prop:bowl-logarithmic-asymptotic} and dividing by $c$ gives the expansion. The speed changes inversely with spatial dilation, as required by parabolic scaling.
\end{proof}
